\documentclass[a4paper,11pt]{article}
\def\email#1{{\tt#1}}

\def\N{\mbox{\rm{I}\hspace{-.15em}\rm{N}}}
\def\Z{\mbox{\rm{Z}\hspace{-.28em}\rm{Z}}}
\def\RR{\mbox{\rm{I}\hspace{-.15em}\rm{R}}}
\def\C{\mbox{\rm{l}\hspace{-.50em}\rm{C}}}
\def\Q{\mbox{\rm{l}\hspace{-.50em}\rm{Q}}}
\def\F{\mbox{\rm{I}\hspace{-.20em}\rm{F}}}

\usepackage{amssymb,amsmath}
\usepackage{url}
\usepackage{fullpage}
\begin{document}

\newtheorem{theo}{Theorem}
\newtheorem{coro}{Corollary}[theo]
\newtheorem{lemme}{Lemma}
\renewcommand{\thecoro}{\thetheo.\arabic{coro}}



\parindent=0cm


\section*{\center IN100 Initiation au Python\\
Contrôle continu 2}

\vspace{1 cm}

\begin{table}[h!]
\centering
\renewcommand{\arraystretch}{1.4} % augmente la hauteur des lignes
\begin{tabular}{|p{10cm}|p{3cm}|p{2cm}|}
\hline
\textbf{Étudiant :} & \textbf{Heure Début }  & \textbf{Note (/5) }  \\
& & \\
\hline
\textbf{Correcteur :} & \textbf{Heure Fin } & \\
& & \\
\hline

\hline
\end{tabular}
\end{table}


\subsection*{Consignes:} Vous avez 20 minutes pour coder la solution. Seules les types et les instructions vues en cours sont autorisées. L'usage d'une fonction 
venant d'une bibliothèque qui n'a pas été vue en cours ne sera pris en compte. 


\subsection*{Sujet:}

\begin{itemize}
\item[1.] Écrire une fonction \texttt{lettres\_identiques}, qui prend en paramètre un mot et qui renvoie \texttt{True} s’il existe deux lettres identiques placées à deux indices consécutifs, \texttt{False} sinon.  
Vous ferez attention aux chaînes vides ou d’un seul caractère.
\vspace*{-0.1 cm}

Par exemple :
\begin{verbatim}
    lettres_identiques("python")     # renvoie False
    lettres_identiques("programmation")  # renvoie True  (les deux 'm')
\end{verbatim}

\item[2.] Écrire une fonction \texttt{couple\_lettres\_identiques}, qui prend en paramètre une chaîne de caractères contenant plusieurs mots et qui renvoie, sous forme de liste de \texttt{Tuple}, tous les couples de lettres identiques placées à deux indices consécutifs.  
Si aucun couple n’est trouvé, la fonction doit renvoyer une liste vide. L'usage des caractères " " consécutifs ne sera pas pris en compte.  

\vspace*{-0.1 cm}

Par exemple :
\begin{verbatim}
    couple_lettres_identiques("belle adresse") # renvoie [('l','l'), ('s','s')]
\end{verbatim}
\end{itemize}

\vspace*{-0.3 cm}


\begin{table}[h!]
\centering
\renewcommand{\arraystretch}{1.4} % augmente la hauteur des lignes
\begin{tabular}{|p{15cm}|}
\hline
\textbf{Remarques étudiant :}    \\
 \\
 
\hline
\textbf{Remarques correcteur :}  \\
 \\
\hline
\end{tabular}
\end{table}

\newpage


\section*{\center IN100 Initiation au Python\\
Contrôle continu 2}


\begin{table}[h!]
\centering
\renewcommand{\arraystretch}{1.4} % augmente la hauteur des lignes
\begin{tabular}{|p{10cm}|p{3cm}|p{2cm}|}
\hline
\textbf{Étudiant :} & \textbf{Heure Début }  & \textbf{Note (/5) }  \\
& & \\
\hline
\textbf{Correcteur :} & \textbf{Heure Fin } & \\
& & \\
\hline

\hline
\end{tabular}
\end{table}

\subsection*{Consignes:} Vous avez 20 minutes pour coder la solution. Seules les types et les instructions vues en cours sont autorisées. Tout usage d'une fonction 
venant d'une bibliothèque qui n'a pas été vue en cours ne sera pris en compte. 


\subsection*{Sujet:}

\begin{itemize}
\item[1.] \'Ecrire une fonction qui prend en argument un mot et renvoie le nombre de voyelles qu'il contient. Si une voyelle apparaît plusieurs fois, on la compte 
autant de fois qu'elle apparaît. 
    \item[2.] Écrire une fonction qui calcule le \textbf{mot contenant le plus de voyelles} dans une chaîne de caractères et le renvoie.
 On suppose que la chaîne ne contient que les lettres de l’alphabet (sans accents) et le caractère \texttt{" "}.

 \vspace*{0.2 cm}
 Par exemple, sur la chaîne suivante :
  
 
\begin{verbatim}
chaine = "tu etais formidable"
\end{verbatim}
    la fonction renverra \texttt{"formidable"}.
% \item[2.] On donne une liste contenant uniquement des chaînes de caractères, chacune représentant une phrase.
% Écrire une fonction qui, à partir de cette liste, calcule le mot contenant le plus de voyelles et renvoie l’indice de la chaîne dans la liste ainsi que le mot.

%Par exemple, sur la liste :
%\vspace*{-0.2 cm}

%\begin{verbatim}chansons = ["tu etais formidable", "Sous le vent resiste belle",  "Alexandrie Alexandra", "Ellе est pas belle la viе"]\end{verbatim}
%votre fonction devra renvoyer \texttt{(3, "Alexandrie")}. 

\end{itemize}
 \vspace*{0.2 cm}

\begin{table}[h!]
\centering
\renewcommand{\arraystretch}{1.4} % augmente la hauteur des lignes
\begin{tabular}{|p{15cm}|}
\hline
\textbf{Remarques étudiant :}    \\
 \\
 \\
\hline
\textbf{Remarques correcteur :}  \\
 \\
 \\
\hline
\end{tabular}
\end{table}


\newpage


\section*{\center IN100 Initiation au Python\\
Contrôle continu 2}


\begin{table}[h!]
\centering
\renewcommand{\arraystretch}{1.4} % augmente la hauteur des lignes
\begin{tabular}{|p{10cm}|p{3cm}|p{2cm}|}
\hline
\textbf{Étudiant :} & \textbf{Heure Début }  & \textbf{Note (/5) }  \\
& & \\
\hline
\textbf{Correcteur :} & \textbf{Heure Fin } & \\
& & \\
\hline

\hline
\end{tabular}
\end{table}
\vspace*{-0.4 cm}
\subsection*{Consignes:} Vous avez 20 minutes pour coder la solution. Seules les types et les instructions vues en cours sont autorisées. Tout usage d'une fonction 
venant d'une bibliothèque qui n'a pas été vue en cours ne sera pris en compte. 

\vspace*{-0.3 cm}
\subsection*{Sujet:}


On souhaite créer un \textbf{chiffrement simple} basé sur l'\textbf{inversion de blocs de lettres}.
Le principe de chiffrement est le suivant: on découpe un mot en \textbf{blocs de taille $n$},  et
on inverse les lettres à l’intérieur de chaque bloc.  Le mot chiffré est obtenu en concaténant tous les blocs inversés.


\subsection*{Exemple}

Mot : \texttt{PYTHON}  avec taille des blocs 2  

\begin{itemize}
    \item Découpage en blocs : \texttt{["PY", "TH", "ON"]}  
    \item Inversion des lettres dans chaque bloc : \texttt{["YP", "HT", "NO"]}  
    \item Mot chiffré : \texttt{YPHTNO}  
\end{itemize}

Le dernier bloc peut être plus court si le mot n’est pas un multiple de la taille du bloc.
\vspace*{-0.3 cm}
\begin{itemize}
\item[1.] Écrire une fonction \texttt{chiffrer\_blocs} qui prend en entrée le mot à chiffrer et la taille du bloc et renvoie le chiffré. 
   

\item[2.] Écrire une fonction \texttt{dechiffrer\_blocs} qui prend en entrée le mot à chiffrer et la taille du bloc et renvoie le mot en clair. 

\end{itemize}

\vspace*{-0.2 cm}
\begin{table}[h!]
\centering
\renewcommand{\arraystretch}{1.4} % augmente la hauteur des lignes
\begin{tabular}{|p{15cm}|}
\hline
\textbf{Remarques étudiant :}    \\
 \\
\hline
\textbf{Remarques correcteur :}  \\
 \\
\hline
\end{tabular}
\end{table}


\newpage

\section*{\center IN100 Initiation au Python\\
Contrôle continu 2}


\begin{table}[h!]
\centering
\renewcommand{\arraystretch}{1.4} % augmente la hauteur des lignes
\begin{tabular}{|p{10cm}|p{3cm}|p{2cm}|}
\hline
\textbf{Étudiant :} & \textbf{Heure Début }  & \textbf{Note (/5) }  \\
& & \\
\hline
\textbf{Correcteur :} & \textbf{Heure Fin } & \\
& & \\
\hline

\hline
\end{tabular}
\end{table}

\vspace*{-0.4 cm}
\subsection*{Consignes:} Vous avez 20 minutes pour coder la solution. Seules les types et les instructions vues en cours sont autorisées. Tout usage d'une fonction 
venant d'une bibliothèque qui n'a pas été vue en cours ne sera pris en compte. 

\vspace*{-0.2 cm}

\subsection*{Sujet:}   

Dans certaines séquences d'ADN, on observe des motifs dits \textbf{palindromiques} :
un fragment est suivi de sa version complémentaire inversée.  
%Ces structures symétriques jouent un rôle important dans la reconnaissance de sites spécifiques par certaines enzymes.
On souhaite modéliser ce phénomène à l’aide de chaînes de caractères représentant des bases d’ADN.
On suppose  que l’ADN est composé de quatre bases :
\vspace*{-0.2 cm}
\[
A, \; T, \; C, \; G
\]
et que les bases complémentaires sont :
\[
A \leftrightarrow T, \quad C \leftrightarrow G.
\]

\begin{itemize}
\item[1.]  Écrire une fonction \texttt{complement} qui prend en entrée  une chaîne d'ADN et renvoie la séquence complémentaire de cette chaîne. 
\vspace*{-0.2 cm}
\begin{verbatim}
   "ATCG" → "TAGC"
\end{verbatim}

\vspace*{-0.2 cm}

\item[2.]  On définit une séquence miroir comme une séquence formée en ajoutant à une séquence donnée
son complément inversé.

\textbf{Exemple :}
\vspace*{-0.2 cm}
\begin{eqnarray*}
S_0 &= &\text{"A"}, \quad S_1 = \text{"A"} + \text{complément inversé de "A"} = \text{"A" + "T"} = \text{"AT"}, \quad \\
S_2 &=& "AT" + \text{complément inversé de "AT"} = \text{"AT" + "AT"} = \text{"ATAT"}
\end{eqnarray*}

Écrire une fonction \texttt{sequence\_miroir} qui à  partir d'une séquence de départ $S_0$ construit la séquence miroir obtenue après \texttt{n} étapes.
\end{itemize}
   
\vspace*{-0.2 cm}

\begin{table}[h!]
\centering
\renewcommand{\arraystretch}{1.4} % augmente la hauteur des lignes
\begin{tabular}{|p{15cm}|}
\hline
\textbf{Remarques étudiant :}    \\
 \\
\hline
\textbf{Remarques correcteur :}  \\
 \\
\hline
\end{tabular}
\end{table}

\newpage

\section*{\center IN100 Initiation au Python\\
Contrôle continu 2}


\begin{table}[h!]
\centering
\renewcommand{\arraystretch}{1.4} % augmente la hauteur des lignes
\begin{tabular}{|p{10cm}|p{3cm}|p{2cm}|}
\hline
\textbf{Étudiant :} & \textbf{Heure Début }  & \textbf{Note (/5) }  \\
& & \\
\hline
\textbf{Correcteur :} & \textbf{Heure Fin } & \\
& & \\
\hline

\hline
\end{tabular}
\end{table}
\vspace{-0.6 cm}
\subsection*{Consignes:} Vous avez 20 minutes pour coder la solution. Seules les types et les instructions vues en cours sont autorisées. Tout usage d'une fonction 
venant d'une bibliothèque qui n'a pas été vue en cours ne sera pris en compte. 

\subsection*{Sujet:}


Vous préparez vos prochaines vacances et hésitez entre plusieurs destinations.
Chaque destination a un coût total estimé. Vous souhaitez choisir une destination qui respecte votre budget.
Vous avez en tête un montant minimum \texttt{budget\_min} et un montant maximum \texttt{budget\_max} que vous souhaitez dépenser.
On dispose d’une liste de prix appelée \texttt{liste\_prix}, contenant le coût total de chaque voyage :

\begin{verbatim}
liste_prix = [450.0, 520.5, 610.0, 670.0, 720.5, 750.0, 810.0, 920.0, 1050.0, 1200.0]
\end{verbatim}


\begin{itemize}
\item[1.] Écrire une fonction \texttt{destinations\_accessibles(liste\_prix, budget\_min, budget\_max)} qui
 prend en argument la liste des prix, un budget minimum et un budget maximum 
et renvoie la liste des destinations dont le coût est compris entre \texttt{budget\_min} et \texttt{budget\_max} inclus.
\item[2.]  Votre ami préfère choisir la destination dont le prix est le plus proche du prix moyen de toutes les destinations proposées.
Écrire une fonction \texttt{destination\_moyenne} qui renvoie qui calcule le prix moyen des voyages et  renvoie le coût de la destination le plus proche de cette moyenne.
\end{itemize}


%\vspace*{-0.3 cm}
\begin{table}[h!]
\centering
\renewcommand{\arraystretch}{1.4} % augmente la hauteur des lignes
\begin{tabular}{|p{15cm}|}
\hline
\textbf{Remarques étudiant :}    \\
 \\
 \\
\hline
\textbf{Remarques correcteur :}  \\
 \\
 \\
\hline
\end{tabular}
\end{table}





\end{document}
